\documentclass{report}

\usepackage{blueprint}

\title{An Elementary Proof of Daboussi's Prime Number Theorem}
\author{Guangxuan Pan (潘广轩)}
\date{\today}

\begin{document}
\maketitle
\tableofcontents

\chapter{Introduction}

\section{The Main Theorem}

The goal of this project is to formalize the elementary proof of Daboussi's Prime Number Theorem, which avoids the use of complex analysis and relies solely on elementary number-theoretic methods.

\begin{theorem}[Daboussi's Prime Number Theorem]\label{thm:daboussi_pnt}
  Let $\pi(x)$ denote the number of primes less than or equal to $x$. Then
  \[
  \lim_{x \to \infty} \frac{\pi(x) \log x}{x} = 1.
  \]
\end{theorem}

\begin{leanok}
  theorem daboussi_pnt : ...
\end{leanok}

\chapter{Key Lemmas}

\section{Selberg's Identity}

The proof crucially relies on Selberg's fundamental identity, which is the cornerstone of elementary proofs of the PNT.

\begin{lemma}[Selberg's Identity]\label{lem:selberg}
  For $x \geq 1$, we have
  \[
  \psi(x) \log x + \sum_{n \leq x} \Lambda(n) \psi\left(\frac{x}{n}\right) = 2x \log x + O(x),
  \]
  where $\psi$ is Chebyshev's psi-function and $\Lambda$ is the von Mangoldt function.
\end{lemma}

\begin{leanok}
  lemma selberg_identity : ...
\end{leanok}

\section{The Chebyshev Functions}

We will need the standard equivalences between different forms of the Prime Number Theorem.

\begin{lemma}[Equivalence of PNT forms]\label{lem:pnt_equivalence}
  The following are equivalent:
  \begin{enumerate}
    \item $\pi(x) \sim \frac{x}{\log x}$,
    \item $\psi(x) \sim x$,
    \item $\vartheta(x) \sim x$.
  \end{enumerate}
\end{lemma}

\begin{leanok}
  lemma pnt_equivalence : ...
\end{leanok}

\chapter{Proof of the Main Theorem}

\section{High-Level Strategy}

The proof proceeds in the following steps:
\begin{enumerate}
  \item Prove Selberg's identity (\Cref{lem:selberg}).
  \item Use the identity to show that $\psi(x) = x + o(x)$.
  \item Conclude the PNT via the equivalence of Chebyshev functions (\Cref{lem:pnt_equivalence}).
\end{enumerate}

\section{From Selberg to PNT}

\begin{proof}[Proof of \Cref{thm:daboussi_pnt}]
  Assume we have established Selberg's identity (\Cref{lem:selberg}).
  We then perform a delicate analysis of the error terms to show that
  \[
  \psi(x) = x + o(x).
  \]
  By \Cref{lem:pnt_equivalence}, this implies
  \[
  \pi(x) \sim \frac{x}{\log x},
  \]
  which is the desired result.
\end{proof}

\end{document}